% SPDX-License-Identifier: Apache-2.0
\section{Feasibility, resources, and the constraints we already know}
\label{sec:feasibility}

\textbf{What the PoC sprint would build.} The same protocol on the issuer's own authorisation
stream: fit the incumbent scorer, freeze four blocks, calibrate the band, certify the
band-conditional false-decline rate at a level the business picks. Every component exists and
has been run end to end here, so the sprint is a port and a re-calibration rather than a build:
three stages of roughly a month with a month of slack --- split integrity, then calibration and
certification, then the quantum re-ranker, which is deliberately last because it is the only
stage that can fail to produce a result.

\textbf{Compute, measured rather than estimated.} The classical core is cheap: the tuned
baseline fits the in-band problem in \ClaimBandGbdtSeconds{} seconds, and conformal calibration
is a sort over the calibration block. The expensive component is the quantum-inspired arm, and
its cost is known --- \ClaimSweepJobs{} full-scale tensor-network fits over
\ClaimSweepTrainRows{} training rows and 431 sites took \ClaimSweepSecondsMin{} to
\ClaimSweepSecondsMax{} seconds each on the one
\Record{RTX~6000 PRO Blackwell (96\,GB), CUDA 13, Python 3.12} workstation every number here was
measured on. Because per-step cost is flat in the bond dimension (Section~\ref{sec:hybrid}),
capacity is nearly free and chain length is the cost, so a resource estimate scales with the
feature count and not with the ansatz. No QPU is required: both arms are simulator-only, which the
statement permits, and the screens ran on \ClaimScreenSamples{} stratified rows across
\ClaimScreenConfigs{} configurations.

\textbf{Latency, against the budget the statement gives.} Of the
\ClaimLatencyAuthorisationBudgetMs{}\,ms authorisation ceiling about \ClaimLatencyNetworkMs{}\,ms
is network, leaving roughly \ClaimLatencyBudgetResidualMs{}\,ms issuer-side. Scoring one
authorisation costs \ClaimLatencyScorerMedian{}\,ms, \ClaimLatencyScorerTail{}\,ms at the tail;
one in-band kernel evaluation would cost \ClaimLatencyKernelMedian{}\,ms ---
\ClaimLatencyKernelVersusScorer{} times as much, its tail taking
\ClaimLatencyKernelBudgetShare\,\% of that budget. The classical core is therefore not the latency
risk and the kernel is: it fits only because the band rations it. Both figures bound model cost
on this workstation, single-threaded, rather than production latency.

\textbf{Data.} The method needs settled labels --- a chargeback window closed before the block
boundary --- and an entity key to cluster on. Both are ordinary in an issuer's data and only
approximated here.

\textbf{Constraints we have measured, and would carry into Phase II.} The population is
conditional on an incumbent: this file contains only transactions an existing stack approved,
because declined attempts cannot generate a chargeback. That is exactly the population a
champion/challenger comparison runs on, so it is the right restriction --- but no result
extrapolates past it. The label is partly an entity flag, since a chargeback propagates to
later transactions sharing an account, email or billing address, and Shapley attribution makes
that concrete: \ClaimAttributionTopShare{} of the in-band model's total contribution falls on
\texttt{card1} alone, across \ClaimAttributionFeatures{} features. A governance review would
read that as a model which mostly recognises cards --- the strongest argument for running this
on an issuer's own data, where an entity key is a key and not a predictor. The certificate is
sample-starved by design: conditioning on a narrow band is what makes the guarantee bind the
component we care about, and also what leaves \ClaimBandLegitCalMin{} legitimate rows to certify
with. A larger calibration block or a wider band buys tighter $\alpha$; the method is unchanged.

\textbf{What would change the quantum conclusion.} An encoding demonstrably outside the RBF
family, at a bandwidth that does not concentrate, would pass the same screens unmodified --- we
did not find one. For the tensor network the honest requirement is a comparison powered for
cross-family noise: resolving an effect the size of the interval in Section~\ref{sec:quantum}
needs about \ClaimMdeSampleFactor{} times the in-band evaluation set, which a wider band or a
longer window would supply. And the negative is about aggregate discrimination, not the
operating point --- the screens read the Gram spectrum and H4 compares average precision across
the band, while H2, which would have tested decisions at the certified threshold, never ran. A
method that does not improve ranking can still improve a decision at a low false-positive
threshold, which is the regime this certificate operates in and which this study did not
measure.
